\makeabstract
    {
    Exploring the Photophysics of Mn2+ Dopant in ZnSe QD System
    }
    {
    Jocelyn MacDonough\textsuperscript{1}, Eunbyeol Gi\textsuperscript{1}, Jacob H. Olshansky\textsuperscript{1}
    }
    {
    \textsuperscript{1} Department of Chemistry, Amherst College, Amherst, MA
    }
    {
    To learn more about the Mn:ZnSe system, we aim to produce a series of Mn:ZnSe QDs with different doping amounts to observe their characteristics and differences. We also want to adjust the synthesis procedure to tune the size and produce uniform dots. We aim to also optimize the EPR settings as Mn has a very strong signal. To further understand the spin of Mn2+, we will do CW-EPR to characterize the spin of Mn2+ at varying concentrations in ZnSe QDs.
    }
    {
    Mn:ZnSe quantum dots are synthesized in an air free environment, with the degassing and heating times modified to tune QD size. UV-Vis spectroscopy and emission are used to observe absorption, emission, and relative quantum yield. Electron paramagnetic resonance (EPR) spectroscopy is performed to observe the spin of Mn2+. To determine size and doping percent of the QDs, transmission electron microscopy and ICP mass spectroscopy are performed. 
    }
    {
    The results of this study suggest that the Mn:ZnSe QDs synthesized express similar characteristics to those in literature. Samples with higher inputs of Mn have lower lifetimes and higher emission values, as expected. The EPR spectra of samples also display the six peaks expected. However, with very high inputs of Mn, samples experience signal disruption, leading to lower emission, or Mn-Mn coupling, leading to a different EPR signal. 
    }
    {
    The observations from this study demonstrate potential in synthesizing a broader series of Mn:ZnSe QDs with different Mn2+ doping percentages. Beyond just characterization, we also aim to determine the electron transfer rate, through methods such as quenching or shelling, and photophysics of Mn:ZnSe QDs.
    }

    \newpage
    